\section{Searching for a closed form}
\label{a:sec:GA}

Before the connection of \cref{a:sec:koenigs} was appreciated, several methods were
tried in the hope of an elementary formula for $A(p)$. The negative result is worth
recording, and so is the manner of the failure.

The first method was the continuum route discussed in \cref{a:sec:asymptotic}:
treat the map as a differential equation in the large-generation limit and solve
it. This gives the shape of the decay but leaves the constant of integration
undetermined, which is to say it leaves $A(p)$ undetermined.

The second was symbolic regression by genetic algorithm. The idea is simple:
generate fifty candidate functions at random, measure how closely each reproduces
the target curve, retain the best five, and build a fresh population of fifty by
varying those; then measure again, and repeat the cycle of variation and selection
many times over \cite{stanley2002evolving,cohen2024physics}. The first
implementation searched a space of polynomials, on the hunch that the answer might
take that form and because such functions are simple to encode: a genome is a
coefficient list $[a_0,\dots,a_k]$ for $a_0+a_1p+\cdots+a_kp^k$ up to some maximum
order, drawn uniformly at random. It became apparent quickly that many
substantially different formulae produce curves indistinguishable from the target
by eye. Supposing that a true formula might have rational coefficients, a second
version searched $\sum_i(a_i/b_i)p^i$ with integer $[a_i],[b_i]$, and a third
extended the space to ratios of two such polynomials. The best that emerged merely
fitted the high-precision numerical data well:
\begin{equation}
\label{a:eq:A1hat}
\hat A_1(p) = \frac{\tfrac12(1-2p)(1+2p)}
{1+\tfrac12p-\tfrac52p^2-\tfrac65p^3+\tfrac65p^4} .
\end{equation}
This passes through $(0,\tfrac12)$ and $(\tfrac12,0)$ and looks convincing, but it
fails the diagnostic of \cref{a:eq:Aasym}: differentiating at $p=\tfrac12$ gives a
gradient of $-2/0.55 = -3.636$ rather than the required $-4$.

A later attempt used a symbolic tree search over progressively nested elementary
functions of all kinds. The conclusion was the same. Many functions of widely
varying form resemble the target closely, and no satisfactory simple formula
emerges even when candidates with fewer components are explicitly favoured. The
expressions returned were typically deep nestings of trigonometric and iterated
exponential terms carrying a dozen or more fitted decimal constants, and they were
not interpretable in any useful sense; representative candidates
$\hat A_2$, $\hat A_3$ and $\hat A_4$, with their plots, are quarantined in
\cref{a:app:catalogue} as records of the fitting exercise rather than proposed
identities.

The recurring defect is the same in every case. The candidates either fail to pass
through $(\tfrac12,0)$ or fail to arrive there with gradient $-4$, and near
criticality this matters enormously: as $A\to0$, an error in $A$ of vanishing
absolute size produces an $O(1)$ --- eventually unbounded --- error in the
quasi-stationary mean $1/A$. A formula can look perfect on a plot of $A$ and be
useless on a plot of $1/A$.

A pragmatic patch is to splice the asymptotic onto a fitted curve, forcing the
correct critical branch; the construction is recorded in \cref{a:app:catalogue},
and it does work. But if one is going to splice, it is better to splice onto
something exact, and \cref{a:sec:practical} does so.

The deeper lesson points directly at the next section. That an enormous number of
very different formulae can behave almost identically over a bounded interval is
not merely an inconvenience of the method; it is a symptom. Curve fitting cannot
distinguish a function that has a closed form from one that does not.
